\section{Selected Organization and Working Environment}

\paragraph{Background and services.}
Google's SRE book says that SREs at Google operate services such as planet-scale storage, email for hundreds of millions of users and web search, which is where Google began \cite{sre-book}. The ``site'' in SRE first meant keeping the google.com website running, and SRE now also runs services that are not websites, from internal infrastructure such as Bigtable to Google Cloud Platform for external developers \cite{sre-book}. The book names Search, Ads, Gmail and Photos as services whose reliability targets Google analyses \cite{sre-book}. SRE itself started when the book's author, who joined in 2003, was tasked with running a ``Production Team'' of seven engineers \cite{sre-book}. The books give no company history, revenue or customer counts beyond these examples.

\paragraph{How SRE is organized.}
SRE and product development are separate organizations at Google, each with its own focus, priorities and management \cite{workbook}. Product teams own their service by default, and SRE is not available or warranted for most services \cite{workbook}. Product teams fund the growth of SRE with new hires when a product succeeds \cite{workbook}. SRE has infrastructural teams, service teams and horizontal product teams \cite{sre-book}. It is not a command-and-control organization, and because of emergency response and pager rotations it is spread over at least a few time zones \cite{sre-book}. Google has hundreds of on-call rotations, and many developer teams, not only SRE, are on call for their own service \cite{workbook}. The Workbook also says the efficacy of SRE at Google is tied to Google's widely accessible unified codebase \cite{workbook}. The books give neither the size of the SRE organization nor its office locations.

\paragraph{Working environment and culture.}
SRE at Google inherited much of its culture from the surrounding company \cite{workbook}. Google says it assumes failure can happen at any time, and that it does not get its culture right all the time \cite{bsrs}. Google prefers multi-site teams in part because night shifts harm people's health \cite{sre-book}, and it compensates on-call work with time off or cash, capped at a proportion of overall salary \cite{workbook}. All new SREs attend SRE EDU, a weeklong course on key concepts, tooling and platforms, and can later spend a week with another team through SRE Exchange \cite{workbook}. The books do not describe offices, working hours, remote-work policy or pay beyond on-call compensation.

\paragraph{Opportunities for students.}
The books mention interns only in two case studies: an SRE team that expected three interns to help its workload \cite{workbook}, and the Chrome security team, which has converted interns to full-time employees \cite{bsrs}. Engineers from product development can spend six months in an SRE team through the Mission Control program \cite{workbook}. We found no description of an SRE internship program, its eligibility, locations or how to apply. That would need Google's careers pages, which we did not use as a source apart from one job advertisement.
